% Part: lambda-calculus
% Chapter: lambda-definability
% Section: lambda-definable-recursive

\documentclass[../../../include/open-logic-section]{subfiles}

\begin{document}

\olfileid{lam}{dfl}{ldr}
\olsection{\usetoken{S}{lambda definable} فنکشن بازگشتی نیں}

سارے جزوی بازگشتی فنکشن !!{lambda definable} نیں، تے اُلٹی گل
وی سچ اے۔ یعنی سارے !!{lambda definable} فنکشن جزوی بازگشتی
نیں۔

\begin{thm}
  \ollabel{thm:lambda-computable} جے جزوی فنکشن $f$،
  !!{lambda definable} اے، تاں اوہ جزوی بازگشتی اے۔
\end{thm}

\begin{proof}
  اسیں صرف ثبوت دا خاکہ دیندے آں۔ پہلاں $\lambd$-ٹرماں نوں
  عددی روپ دیندے آں، یعنی $\lambd$-ٹرماں نوں قاعدے نال گوڈل
  نمبر دیندے آں۔ ایس لئی سلسلیاں دا عام کوڈ ورتدے آں جیہڑا
  اولی عدداں دیاں گھاتاں نال بندا اے۔ فیر اک جزوی بازگشتی
  فنکشن $\fn{normalize}(t)$ متعین کردے آں جیہڑا دلیل وجوں
  لامبڈا ٹرم دا گوڈل نمبر~$t$ لیندا اے۔ جے ٹرم دا نارمل روپ
  ہووے، تاں اوہدا گوڈل نمبر واپس دیندا اے۔ ہور صورت وچ
  غیر معین اے۔ فیر دو جزوی بازگشتی فنکشن $\fn{toChurch}$
  تے $\fn{fromChurch}$ متعین کردے آں جیہڑے قدرتی عدداں
  نوں اوہناں دے چرچ عدداں دے گوڈل نمبراں وچ تے واپس بدلدے نیں۔

  ایہناں بازگشتی فنکشوناں نال اسیں فنکشن~$f$ نوں جزوی بازگشتی
  فنکشن وجوں متعین کر سکدے آں۔ کوئی $\lambd$-ٹرم~$F$ اے
  جیہڑا~$f$ نوں !!{lambda define}s۔ $f(n_1, \dots, n_k)$ دا
  حساب کرن لئی پہلاں $\fn{toChurch}(n_i)$ نال متعلقہ چرچ
  عدداں دے گوڈل نمبر حاصل کرو۔ ایہناں نوں $\Gn{F}$ نال
  جوڑ کے ٹرم $F \num{n_1}\dots\num{n_k}$ دا گوڈل نمبر بناؤ۔
  ہُن ایس گوڈل نمبر اُتے $\fn{normalize}$ لاؤ۔ جے
  $f(n_1, \dots, n_k)$ معین اے، تاں $F \num{n_1}\dots\num{n_k}$
  دا نارمل روپ اے تے اوہ لازماً چرچ عدد اے۔ ہور صورت وچ
  اوہدا کوئی نارمل روپ نہیں، تے ایس لئی
  \[\fn{normalize}(\Gn{F\num{n_1}\dots\num{n_k}})\]
  غیر معین اے۔ آخر وچ نارمل کیتے ٹرم دے گوڈل نمبر اُتے
  $\fn{fromChurch}$ لاؤ۔
\end{proof}

\end{document}
